% Options for packages loaded elsewhere
\PassOptionsToPackage{unicode}{hyperref}
\PassOptionsToPackage{hyphens}{url}
\documentclass[
]{article}
\usepackage{xcolor}
\usepackage{amsmath,amssymb}
\setcounter{secnumdepth}{-\maxdimen} % remove section numbering
\usepackage{iftex}
\ifPDFTeX
  \usepackage[T1]{fontenc}
  \usepackage[utf8]{inputenc}
  \usepackage{textcomp} % provide euro and other symbols
\else % if luatex or xetex
  \usepackage{unicode-math} % this also loads fontspec
  \defaultfontfeatures{Scale=MatchLowercase}
  \defaultfontfeatures[\rmfamily]{Ligatures=TeX,Scale=1}
\fi
\usepackage{lmodern}
\ifPDFTeX\else
  % xetex/luatex font selection
\fi
% Use upquote if available, for straight quotes in verbatim environments
\IfFileExists{upquote.sty}{\usepackage{upquote}}{}
\IfFileExists{microtype.sty}{% use microtype if available
  \usepackage[]{microtype}
  \UseMicrotypeSet[protrusion]{basicmath} % disable protrusion for tt fonts
}{}
\makeatletter
\@ifundefined{KOMAClassName}{% if non-KOMA class
  \IfFileExists{parskip.sty}{%
    \usepackage{parskip}
  }{% else
    \setlength{\parindent}{0pt}
    \setlength{\parskip}{6pt plus 2pt minus 1pt}}
}{% if KOMA class
  \KOMAoptions{parskip=half}}
\makeatother
\usepackage{longtable,booktabs,array}
\usepackage{caption}
\captionsetup[table]{skip=6pt}
\newcounter{none} % for unnumbered tables
\usepackage{calc} % for calculating minipage widths
% Correct order of tables after \paragraph or \subparagraph
\usepackage{etoolbox}
\makeatletter
\patchcmd\longtable{\par}{\if@noskipsec\mbox{}\fi\par}{}{}
\makeatother
% Allow footnotes in longtable head/foot
\IfFileExists{footnotehyper.sty}{\usepackage{footnotehyper}}{\usepackage{footnote}}
\makesavenoteenv{longtable}
\setlength{\emergencystretch}{3em} % prevent overfull lines
\providecommand{\tightlist}{%
  \setlength{\itemsep}{0pt}\setlength{\parskip}{0pt}}
\usepackage{bookmark}
\IfFileExists{xurl.sty}{\usepackage{xurl}}{} % add URL line breaks if available
\urlstyle{same}
% fallback for those not using the hyperref driver hyperxmp:
\makeatletter
\@ifundefined{xmpquote}{\newcommand{\xmpquote}[1]{#1}}{}
\makeatother
\hypersetup{
  hidelinks,
  pdfcreator={LaTeX via pandoc}}

\author{}
\date{}

\begin{document}

\section{S\&P 500 Quantitative Data Analysis --- Master Glossary \&
Lexicon (Grand Lexique
Bilingue)}\label{sp-500-quantitative-data-analysis-master-glossary-lexicon-grand-lexique-bilingue}

\begin{quote}
\textbf{Course}: Quantitative Data Analysis (M1 -- S7, Course code:
\texttt{2627\_ECO\_2\_EN\_009}) \textbf{Institution}: EM Normandie
Business School · Programme Grande École \textbf{Lecturer}: Dr.~NGUYEN
Anh-Tuan \textbf{Purpose}: A comprehensive, beginner-friendly bilingual
guide (English concepts with detailed French explanations) covering all
Financial, Data, Statistical, Econometric, ESG, and Jamovi terminology
across all 4 sessions of the project.
\end{quote}

\begin{center}\rule{0.5\linewidth}{0.5pt}\end{center}

\subsection{Table of Contents}\label{table-of-contents}

\begin{enumerate}
\def\labelenumi{\arabic{enumi}.}
\tightlist
\item
  \hyperref[1-data-fundamentals--variable-typology]{Data Fundamentals \&
  Variable Typology}
\item
  \hyperref[2-financial--corporate-accounting-terminology]{Financial \&
  Corporate Accounting Terminology}
\item
  \hyperref[3-esg--corporate-governance-iss-framework]{ESG \& Corporate
  Governance (ISS Framework)}
\item
  \hyperref[4-session-1-univariate-statistics--shape-parameters]{Session
  1: Univariate Statistics \& Shape Parameters}
\item
  \hyperref[5-session-2-bivariate-associations-chi-square--correlations]{Session
  2: Bivariate Associations, Chi-Square \& Correlations}
\item
  \hyperref[6-session-3-group-comparisons-welchs-anova--games-howell]{Session
  3: Group Comparisons, Welch's ANOVA \& Games-Howell}
\item
  \hyperref[7-session-4-multiple-linear-regression--econometric-diagnostics]{Session
  4: Multiple Linear Regression \& Econometric Diagnostics}
\item
  \hyperref[8-jamovi-software-operations--costly-student-traps]{Jamovi
  Software Operations \& Costly Student Traps}
\end{enumerate}

\begin{center}\rule{0.5\linewidth}{0.5pt}\end{center}

\subsection{1. Data Fundamentals \& Variable
Typology}\label{data-fundamentals-variable-typology}

\subsubsection{Unit of Observation (Unité
d'observation)}\label{unit-of-observation-unituxe9-dobservation}

\begin{itemize}
\tightlist
\item
  \textbf{Definition}: Exactly what a single row represents in your
  database.
\item
  \textbf{In our project}: \textbf{1 row = 1 publicly traded S\&P 500
  company}.
\item
  \textbf{Why it matters (Slide 11)}: If you cannot finish the sentence
  \emph{``One row = one\ldots{}''}, you do not understand your dataset.
  Having 1 row = 1 firm ensures independence of observations for ANOVA
  and regression.
\end{itemize}

\subsubsection{\texorpdfstring{Sample Size (\(N\) / Observation
count)}{Sample Size (N / Observation count)}}\label{sample-size-n-observation-count}

\begin{itemize}
\tightlist
\item
  \textbf{Definition}: The total number of valid entities observed in
  the study.
\item
  \textbf{In our project}: \(N = 503\) companies (exceeds the course
  minimum requirement of \(N \ge 100\)).
\end{itemize}

\subsubsection{Qualitative / Categorical Variable (Variable
qualitative)}\label{qualitative-categorical-variable-variable-qualitative}

Variables that express names, categories, or labels rather than numeric
quantities. You can count their frequency, but you \textbf{cannot}
compute a mathematical average.

\begin{itemize}
\tightlist
\item
  \textbf{Nominal Variable (Variable nominale)}:

  \begin{itemize}
  \tightlist
  \item
    Categories with \textbf{no natural order or ranking}.
  \item
    \emph{Examples}: Industry sector (\texttt{Sector}), country of
    listing, legal status.
  \item
    \emph{Course rule (Slide 29)}: Must have between \textbf{2 and 5
    categories}. Our project has exactly 5: \emph{Tech \& Comms,
    Healthcare, Finance, Industrials \& Energy, Consumer}.
  \item
    \emph{Golden Rule (Slide 19)}: \textbf{Never compute a mean for a
    nominal variable!} Saying \emph{``Average sector = 2.4''} is
    completely meaningless and results in lost marks.
  \end{itemize}
\item
  \textbf{Ordinal Variable (Variable ordinale)}:

  \begin{itemize}
  \tightlist
  \item
    Categories that possess a \textbf{clear, natural ranking or order},
    but where the distance (gap) between categories is \textbf{not
    equal} or mathematically measurable.
  \item
    \emph{Examples}: Credit rating (AAA \textgreater{} AA \textgreater{}
    A), Customer satisfaction (1 star to 5 stars), Governance tier
    (\emph{Low \textless{} Medium \textless{} High}), Size quartile
    (\emph{Q1 \textless{} Q2 \textless{} Q3 \textless{} Q4}).
  \item
    \emph{Golden Rule (Slide 13)}: The distance between A and B is not
    the same as between B and C. Central tendency is captured by the
    \textbf{mode} or the \textbf{median category}.
  \end{itemize}
\end{itemize}

\subsubsection{Quantitative / Numerical Variable (Variable
quantitative)}\label{quantitative-numerical-variable-variable-quantitative}

Variables that represent measurable or countable numerical values where
arithmetic operations (addition, subtraction, averaging) are valid.

\begin{itemize}
\tightlist
\item
  \textbf{Discrete Variable (Variable discrète)}: Countable integer
  values with no decimals (e.g., Number of employees \texttt{Headcount},
  number of factories, number of board members).
\item
  \textbf{Continuous Variable (Variable continue)}: Measured on an
  unbroken numerical continuum, allowing decimals (e.g.,
  \texttt{Profit\_Margin}, \texttt{Total\_Revenue\_B},
  \texttt{Market\_Cap\_B}, \texttt{Beta}).
\item
  \textbf{Scale Types: Interval vs.~Ratio (Échelle d'intervalle
  vs.~Échelle de rapport)}:

  \begin{itemize}
  \tightlist
  \item
    \emph{Interval scale}: Has an arbitrary zero point (e.g.,
    Temperature in °C: 0°C does not mean ``absence of temperature'', so
    20°C is not ``twice as hot'' as 10°C).
  \item
    \emph{Ratio scale}: Has a \textbf{true, absolute zero} indicating
    complete absence of the quantity (e.g., Turnover = \$0 means zero
    sales; \$20M is truly twice as much as \$10M).
  \item
    \emph{In our project}: \texttt{Total\_Revenue\_B},
    \texttt{Market\_Cap\_B}, and \texttt{Profit\_Margin} are ratio-scale
    continuous variables, enabling full correlation and regression
    modeling.
  \end{itemize}
\end{itemize}

\subsubsection{Dataset Structure: Cross-Sectional vs.~Time-Series
vs.~Panel}\label{dataset-structure-cross-sectional-vs.-time-series-vs.-panel}

\begin{itemize}
\tightlist
\item
  \textbf{Cross-Sectional Dataset (Coupe transversale)}: Multiple
  entities (503 firms) observed at a \textbf{single point in time}
  (e.g., FY2025). This is our project structure.
\item
  \textbf{Time-Series (Série temporelle)}: A single entity observed over
  multiple consecutive periods (e.g., Apple's stock price daily from
  2020 to 2026).
\item
  \textbf{Panel Data (Données de panel)}: Multiple entities observed
  across multiple time periods (\(1\text{ row} = 1\text{ firm-year}\)).
  \emph{Caution}: Panel data introduces autocorrelation, which violates
  standard ANOVA and OLS regression assumptions taught in this course.
\end{itemize}

\subsubsection{Outliers: Detection \& Handling (Valeurs aberrantes /
extrêmes)}\label{outliers-detection-handling-valeurs-aberrantes-extruxeames}

Data points that deviate substantially from the overall pattern of the
distribution. Outliers can distort arithmetic means, inflate standard
deviations, and invalidate normality:

\begin{itemize}
\tightlist
\item
  \textbf{Tukey's \(1.5 \times IQR\) Rule (Règle des moustaches de
  Tukey)}:

  \begin{itemize}
  \tightlist
  \item
    Any observation below the lower fence:
    \(\text{LF} = Q_1 - 1.5 \times IQR\)
  \item
    Any observation above the upper fence:
    \(\text{UF} = Q_3 + 1.5 \times IQR\)
  \item
    Points outside this range appear as individual dots outside the
    boxplot whiskers (e.g., NVIDIA or Apple's extreme market cap; highly
    negative profit margins in turnaround firms).
  \end{itemize}
\item
  \textbf{Z-Score Method (Seuil en écarts-types)}:

  \begin{itemize}
  \tightlist
  \item
    Standardized distance: \(z_i = \frac{x_i - \bar{x}}{s}\). An
    observation with \(|z_i| > 3.0\) is traditionally flagged as an
    extreme outlier in normally distributed data.
  \end{itemize}
\item
  \textbf{Management Strategy (Slide 21)}: Never delete legitimate
  business outliers (like mega-cap firms) without justification!
  Instead, document them, apply non-parametric/robust tests (Median,
  IQR, Welch ANOVA), or use logarithmic transformations.
\end{itemize}

\subsubsection{Missing Data: Mechanisms \& Treatment (Données
manquantes)}\label{missing-data-mechanisms-treatment-donnuxe9es-manquantes}

\begin{itemize}
\tightlist
\item
  \textbf{Missing Completely at Random (MCAR --- Manquant complètement
  au hasard)}: The probability of missingness is unrelated to any
  observed or unobserved variable (e.g., accidental scraping network
  drop on a random ticker).
\item
  \textbf{Missing at Random (MAR --- Manquant au hasard conditionnel)}:
  Missingness depends on other observed variables (e.g., smaller firms
  failing to disclose certain governance sub-scores).
\item
  \textbf{Missing Not at Random (MNAR --- Manquant non au hasard)}:
  Missingness depends on the unobserved value itself (e.g., firms with
  catastrophic ROE having negative equity, making ROE mathematically
  undefined).
\item
  \textbf{Treatment Strategies}:

  \begin{itemize}
  \tightlist
  \item
    \textbf{Listwise Deletion (Suppression par liste / Complete Case
    Analysis)}: Dropping any row containing a missing value across the
    analyzed variables. Safe when MCAR and sample size is large
    (\(N = 503\)).
  \item
    \textbf{Pairwise Deletion (Suppression par paire)}: Used in
    bivariate correlation matrices (calculating \(r\) between variables
    \(X\) and \(Y\) using all rows where both \(X\) and \(Y\) are valid,
    preserving sample size per pair).
  \item
    \textbf{Data Imputation (Imputation de données)}: Replacing missing
    values with mean/median or predictive regression (avoided here to
    prevent artificial variance deflation).
  \end{itemize}
\end{itemize}

\begin{center}\rule{0.5\linewidth}{0.5pt}\end{center}

\subsection{2. Financial \& Corporate Accounting
Terminology}\label{financial-corporate-accounting-terminology}

\subsubsection{S\&P 500 Index (Standard \& Poor's
500)}\label{sp-500-index-standard-poors-500}

A market-capitalization-weighted equity index tracking the 500 largest
public companies listed on stock exchanges in the United States (NYSE
and NASDAQ), representing approximately 80\% of total U.S. equity market
capitalization.

\subsubsection{Ticker Symbol (Symbole
boursier)}\label{ticker-symbol-symbole-boursier}

A unique arrangement of 1 to 5 letters identifying a specific publicly
traded security on an exchange (e.g., \texttt{AAPL} for Apple,
\texttt{MSFT} for Microsoft, \texttt{NVDA} for Nvidia, \texttt{FI} for
Fiserv). Serves as the primary \texttt{ID} variable in our dataset.

\subsubsection{\texorpdfstring{Market Capitalization
(\texttt{Market\_Cap\_B} --- Capitalisation
boursière)}{Market Capitalization (Market\_Cap\_B --- Capitalisation boursière)}}\label{market-capitalization-market_cap_b-capitalisation-boursiuxe8re}

\[\text{Market Cap} = \text{Current Share Price} \times \text{Total Shares Outstanding}\]
The total equity market value of a public company, expressed in
\textbf{Billions of USD (\$B)} in our dataset.

\begin{itemize}
\tightlist
\item
  \emph{Market Cap Quartiles (\texttt{Market\_Cap\_Quartile})}: The S\&P
  500 divided into 4 equal tiers of firm size:

  \begin{itemize}
  \tightlist
  \item
    \textbf{Q1}: Smallest quartile in the index (\$6.5B to \$22.2B)
  \item
    \textbf{Q2}: Mid-sized large-cap (\$22.2B to \$44.5B)
  \item
    \textbf{Q3}: Upper large-cap (\$44.5B to \$94.0B)
  \item
    \textbf{Q4}: Mega-cap giants (\$94.0B to \$5,367.1B, including
    Apple, Microsoft, Nvidia)
  \end{itemize}
\end{itemize}

\subsubsection{\texorpdfstring{Total Revenue (\texttt{Total\_Revenue\_B}
--- Chiffre d'affaires /
Turnover)}{Total Revenue (Total\_Revenue\_B --- Chiffre d'affaires / Turnover)}}\label{total-revenue-total_revenue_b-chiffre-daffaires-turnover}

The total gross amount of income generated by the sale of goods or
services related to the company's primary operations before deducting
any expenses, taxes, or depreciation. Reported in \textbf{Billions of
USD (\$B)}.

\subsubsection{Fiscal Year (FY --- Exercice fiscal) vs.~Calendar Year
(Année
calendaire)}\label{fiscal-year-fy-exercice-fiscal-vs.-calendar-year-annuxe9e-calendaire}

\begin{itemize}
\tightlist
\item
  \emph{Calendar Year}: Strictly January 1 to December 31.
\item
  \emph{Fiscal Year (FY2025)}: The official 12-month accounting period
  used by a corporation to prepare audited financial reports (SEC Form
  10-K).
\item
  \emph{Why it matters}: Many US companies do not close their books on
  December 31. Apple closes in September; Nvidia and Walmart close in
  January; Microsoft closes in June. Using \textbf{FY2025} ensures we
  analyze 12 full months of audited operating figures for every single
  corporation.
\end{itemize}

\subsubsection{Trailing Twelve Months (TTM --- Douze derniers mois
glissants)}\label{trailing-twelve-months-ttm-douze-derniers-mois-glissants}

A financial reporting measurement window that calculates performance
over the most recent 12-month period (the last 4 quarters) regardless of
the fiscal year-end date.

\subsubsection{\texorpdfstring{Net Profit Margin
(\texttt{Profit\_Margin} --- Marge bénéficiaire
nette)}{Net Profit Margin (Profit\_Margin --- Marge bénéficiaire nette)}}\label{net-profit-margin-profit_margin-marge-buxe9nuxe9ficiaire-nette}

\[\text{Profit Margin} = \frac{\text{Net Income}}{\text{Total Revenue}}\]
The percentage of revenue that remains as bottom-line profit after
paying all operating expenses, interest, and taxes. Expressed as a
decimal in our dataset (e.g., \(0.15 = 15\%\)).

\begin{itemize}
\tightlist
\item
  \emph{Advantage in Data Analysis}: Revenue is positive for all
  operating firms, meaning \texttt{Profit\_Margin} is 100\% defined and
  available for all 503 firms without accounting distortion.
\end{itemize}

\subsubsection{\texorpdfstring{Return on Equity (\texttt{ROE} ---
Rentabilité des capitaux
propres)}{Return on Equity (ROE --- Rentabilité des capitaux propres)}}\label{return-on-equity-roe-rentabilituxe9-des-capitaux-propres}

\[\text{ROE} = \frac{\text{Net Income}}{\text{Shareholders' Equity}}\]
Measures how efficiently a company's management uses shareholders'
capital to generate net profits. Expressed as a decimal (e.g.,
\(0.20 = 20\%\)).

\subsubsection{Negative Shareholders' Equity \& Missing ROE (Capitaux
propres
négatifs)}\label{negative-shareholders-equity-missing-roe-capitaux-propres-nuxe9gatifs}

\begin{itemize}
\tightlist
\item
  \emph{The Phenomenon}: In our dataset, 32 prominent corporations
  (AbbVie, Altria, AutoZone, Booking Holdings, Domino's Pizza, etc.)
  show missing/NaN values for ROE.
\item
  \emph{The Financial Reason}: These mature, highly cash-generative
  firms engaged in aggressive, debt-funded \textbf{share buybacks
  (rachats d'actions)} over 10+ years. Buying back stock reduces the
  accounting book equity on the balance sheet. When cumulative buybacks
  exceed retained earnings, book equity becomes negative.
\item
  \emph{Econometric Implication}: Net income divided by negative equity
  produces a mathematically meaningless or negative number. Financial
  databases (Yahoo Finance, Bloomberg) set ROE to \texttt{NaN}. Rather
  than deleting these premier companies (which would introduce survival
  and industry bias), we document this in the \emph{Data Cleaning Log}
  and utilize \texttt{Profit\_Margin} as our universal profitability
  indicator.
\end{itemize}

\subsubsection{\texorpdfstring{Market Beta (\texttt{Beta} --- Volatilité
systématique)}{Market Beta (Beta --- Volatilité systématique)}}\label{market-beta-beta-volatilituxe9-systuxe9matique}

A quantitative measure of a stock's volatility (systematic market risk)
in comparison to the overall market (S\&P 500 index), calculated from a
\textbf{5-year window of monthly returns}:

\begin{itemize}
\tightlist
\item
  \(\beta = 1.0\): Stock moves in tandem with the market.
\item
  \(\beta > 1.0\): Aggressive / high-volatility stock (e.g., high-growth
  tech).
\item
  \(\beta < 1.0\): Defensive / low-volatility stock (e.g., consumer
  staples, utilities).
\item
  \(\beta < 0\): Inverted movement (extremely rare in equities).
\end{itemize}

\subsubsection{Form 10-K (Rapport annuel
SEC)}\label{form-10-k-rapport-annuel-sec}

The comprehensive annual regulatory filing required by the U.S.
Securities and Exchange Commission (SEC) providing a detailed, audited
breakdown of a public company's financial performance.

\subsubsection{GICS (Global Industry Classification Standard ---
Classification
sectorielle)}\label{gics-global-industry-classification-standard-classification-sectorielle}

An industry taxonomy developed in 1999 by MSCI and S\&P Dow Jones
Indices. Classifies all public companies into \textbf{11 primary
sectors}: \emph{Information Technology, Communication Services,
Financials, Health Care, Consumer Discretionary, Consumer Staples,
Industrials, Energy, Utilities, Real Estate, Materials}.

\begin{itemize}
\tightlist
\item
  \textbf{Course Mapping (Slide 29)}: The course strictly limits the
  nominal grouping factor to 2--5 categories. We consolidated the 11 raw
  GICS sectors into \textbf{5 economically coherent groups}: \emph{Tech
  \& Comms, Healthcare, Finance, Industrials \& Energy, Consumer}.
\end{itemize}

\subsubsection{The Fundamental Accounting Equation (Équation
fondamentale du
bilan)}\label{the-fundamental-accounting-equation-uxe9quation-fondamentale-du-bilan}

\[\text{Assets} = \text{Liabilities} + \text{Shareholders' Equity}\]
\[\text{Actif} = \text{Dettes} + \text{Capitaux Propres}\] A company's
economic resources (Assets) are financed either by borrowing from
creditors (Liabilities) or by capital provided and accumulated by owners
(Equity). When liabilities exceed assets (often due to massive
debt-funded share repurchases), accounting equity turns negative.

\subsubsection{The P\&L Income Statement Waterfall (Cascade du compte de
résultat)}\label{the-pl-income-statement-waterfall-cascade-du-compte-de-ruxe9sultat}

The sequential deduction of operating and non-operating costs to arrive
at net bottom-line earnings:

\begin{enumerate}
\def\labelenumi{\arabic{enumi}.}
\tightlist
\item
  \textbf{Gross Revenue (Chiffre d'affaires)}: Total operational sales
  inflow (\texttt{Total\_Revenue\_B}).
\item
  \textbf{Gross Profit (Marge brute)}: Revenue minus Cost of Goods Sold
  (COGS).
\item
  \textbf{EBITDA (Excédent Brut d'Exploitation)}: Earnings Before
  Interest, Taxes, Depreciation, and Amortization. A pure measure of
  operational cash generation.
\item
  \textbf{EBIT / Operating Income (Résultat d'exploitation)}:
  Operational profit after depreciation of physical plant and equipment.
\item
  \textbf{EBT (Résultat avant impôts)}: Operating income minus interest
  on debt.
\item
  \textbf{Net Income (Résultat net)}: Final accounting earnings
  distributed to shareholders after corporate income tax. Divided by
  revenue, this yields our key outcome variable:
  \texttt{Profit\_Margin}.
\end{enumerate}

\subsubsection{Systematic Risk vs.~Idiosyncratic Risk (CAPM /
MEDAF)}\label{systematic-risk-vs.-idiosyncratic-risk-capm-medaf}

Under the Capital Asset Pricing Model (CAPM):

\begin{itemize}
\tightlist
\item
  \textbf{Systematic Market Risk (Risque systématique --- \(\beta\))}:
  Macroeconomic risks affecting the entire market (interest rates,
  inflation, recessions, geopolitical crises). Cannot be eliminated by
  diversification.
\item
  \textbf{Idiosyncratic / Unsystematic Risk (Risque spécifique /
  idiosyncratique)}: Firm-specific risks (e.g., product failure,
  executive fraud, board scandal). Completely diversifiable in a large
  portfolio.
\item
  \textbf{Our Empirical Finding}: Poor corporate governance is often
  assumed to be purely idiosyncratic, but our Session 2 correlation
  proves that weak governance significantly inflates \textbf{systematic
  market sensitivity} (\(r = +0.268, p < .001^{***}\))!
\end{itemize}

\subsubsection{Earnings Per Share (EPS / BPA) \& Financial
Engineering}\label{earnings-per-share-eps-bpa-financial-engineering}

\[\text{EPS} = \frac{\text{Net Income}}{\text{Total Shares Outstanding}}\]
Because executive bonuses and stock option vestings are heavily tied to
EPS targets, corporate executives often use low-interest corporate debt
to buy back company shares from the open market. This reduces the
denominator (Shares Outstanding), artificially boosting EPS even when
top-line business growth is stagnant.

\subsubsection{Weighted Average Cost of Capital (WACC /
CMPC)}\label{weighted-average-cost-of-capital-wacc-cmpc}

The average after-tax rate of return a company must pay to all its
security holders (debt-holders and equity-holders) to finance its
assets. Poor corporate governance increases perceived risk among lenders
and investors, resulting in a \textbf{governance risk premium} that
elevates the firm's WACC and reduces valuation.

\begin{center}\rule{0.5\linewidth}{0.5pt}\end{center}

\subsection{3. ESG \& Corporate Governance (ISS
Framework)}\label{esg-corporate-governance-iss-framework}

\subsubsection{ESG (Environmental, Social, and
Governance)}\label{esg-environmental-social-and-governance}

A tri-pillar framework used by institutional investors to evaluate
corporate sustainability, societal impact, and ethical leadership
alongside traditional financial statements.

\subsubsection{ISS (Institutional Shareholder
Services)}\label{iss-institutional-shareholder-services}

The global leader in proxy voting advisory and corporate governance
research. ISS evaluates publicly listed firms and assigns quantitative
risk ratings based on regulatory disclosures.

\subsubsection{\texorpdfstring{ISS Governance QualityScore
(\texttt{Overall\_Governance\_Risk})}{ISS Governance QualityScore (Overall\_Governance\_Risk)}}\label{iss-governance-qualityscore-overall_governance_risk}

A composite decile score measuring corporate governance risk across 4
core pillars on a \textbf{1 to 10 scale}:

\begin{itemize}
\tightlist
\item
  \textbf{Audit Risk (\texttt{Audit\_Risk})}: Financial integrity,
  accounting oversight, auditor independence, and financial
  restatements.
\item
  \textbf{Board Structure Risk (\texttt{Board\_Risk})}: Independence of
  directors, board diversity, committee structures, and separation of
  CEO and Chairman roles.
\item
  \textbf{Compensation Risk (\texttt{Compensation\_Risk})}: Alignment
  between executive pay and shareholder returns, clawback policies, and
  golden parachutes.
\item
  \textbf{Shareholder Rights Risk (\texttt{Shareholder\_Rights\_Risk})}:
  Voting rights, poison pills, classified boards, and takeover defense
  mechanisms.
\end{itemize}

\begin{quote}
{[}!CAUTION{]} \textbf{Crucial Inversion}: In the ISS methodology,
\textbf{1 represents the lowest risk (best governance practices)}, while
\textbf{10 represents the highest risk (poorest governance practices)}!
A negative correlation between Governance Risk and Profit Margin
therefore implies that \emph{better governance (lower risk score) is
associated with higher profitability}!
\end{quote}

\subsubsection{\texorpdfstring{Governance Risk Tier
(\texttt{Governance\_Risk\_Level} --- Variable
ordinale)}{Governance Risk Tier (Governance\_Risk\_Level --- Variable ordinale)}}\label{governance-risk-tier-governance_risk_level-variable-ordinale}

To enable categorical contingency tables (Chi-Square) and group tests,
the 1--10 continuous score is classified into three ordered tiers:

\begin{itemize}
\tightlist
\item
  \textbf{Low Risk} (Scores 1 to 3): Premier governance compliance
  (\(n = 150\), \(30.2\%\)).
\item
  \textbf{Medium Risk} (Scores 4 to 6): Standard market governance
  (\(n = 149\), \(30.0\%\)).
\item
  \textbf{High Risk} (Scores 7 to 10): Elevated governance friction
  (\(n = 197\), \(39.7\%\), Modal tier).
\end{itemize}

\subsubsection{Agency Theory \& Corporate Governance (Théorie de
l'agence)}\label{agency-theory-corporate-governance-thuxe9orie-de-lagence}

The foundational financial theory of corporate governance (Jensen \&
Meckling, 1976; Fama \& Jensen, 1983):

\begin{itemize}
\tightlist
\item
  \textbf{The Core Conflict}: The \textbf{separation of ownership}
  (principals = shareholders who provide capital) from
  \textbf{operational control} (agents = managers and CEOs who run the
  day-to-day business).
\item
  \textbf{Agency Costs (Coûts d'agence)}: Self-interested managers may
  maximize their own private benefits (excessive pay, corporate jets,
  empire building through unprofitable mergers) rather than maximizing
  long-term shareholder value.
\item
  \textbf{The Role of Governance}: Board oversight, independent
  auditing, executive pay alignment, and shareholder voting rights exist
  precisely to minimize agency friction and ensure managers work in the
  interest of owners.
\end{itemize}

\subsubsection{Key Governance Mechanisms Monitored by
ISS}\label{key-governance-mechanisms-monitored-by-iss}

\begin{itemize}
\tightlist
\item
  \textbf{CEO Duality (Cumul des mandats de PDG et Président du
  Conseil)}: When the Chief Executive Officer also serves as Chairman of
  the Board of Directors. ISS considers this a major conflict of
  interest because the board's primary role is to evaluate and supervise
  the CEO. Having an \textbf{Independent Lead Director} mitigates this
  risk.
\item
  \textbf{Say-on-Pay (Vote consultatif sur la rémunération)}: A
  non-binding shareholder vote mandated under the Dodd-Frank Act
  allowing investors to approve or reject executive compensation
  packages. Repeated low approval (\textless{} 70\%) triggers elevated
  ISS Compensation Risk.
\item
  \textbf{Clawback Provision (Clause de restitution des bonus)}: A
  contractual clause enabling the board to claw back previously awarded
  executive bonuses in the event of financial restatements, fraud, or
  material misconduct.
\item
  \textbf{Poison Pill / Shareholder Rights Plan (Pilule empoisonnée)}: A
  defense mechanism allowing existing shareholders to purchase newly
  issued shares at a steep discount during an uninvited hostile takeover
  bid, severely diluting the hostile acquirer. ISS heavily penalizes
  companies adopting poison pills without shareholder approval.
\item
  \textbf{Classified / Staggered Board (Conseil d'administration
  échelonné)}: A board structure where only a fraction (typically
  one-third) of directors stand for election each year. Prevents an
  activist shareholder or acquirer from replacing the entire board in a
  single proxy contest.
\end{itemize}

\subsubsection{Proxy Advisory \& The ``Big Three'' Asset
Managers}\label{proxy-advisory-the-big-three-asset-managers}

\begin{itemize}
\tightlist
\item
  \textbf{Proxy Advisors (Agences de conseil en vote)}: Firms like ISS
  and Glass Lewis that research shareholder proposals, analyse corporate
  filings, and issue vote recommendations for annual general meetings
  (AGMs).
\item
  \textbf{The Big Three (BlackRock, Vanguard, State Street)}: The
  largest passive index asset managers in the world, holding on average
  over 20\% of the voting shares of S\&P 500 corporations. Because they
  manage passive index funds and cannot simply ``sell'' the stock of a
  poorly run company, they rely heavily on \textbf{ISS governance scores
  and proxy voting} to force corporate governance reform.
\end{itemize}

\begin{center}\rule{0.5\linewidth}{0.5pt}\end{center}

\subsection{4. Session 1: Univariate Statistics \& Shape
Parameters}\label{session-1-univariate-statistics-shape-parameters}

\subsubsection{Univariate Analysis (Analyse
univariée)}\label{univariate-analysis-analyse-univariuxe9e}

The statistical examination of \textbf{one single variable in
isolation}, without analyzing causes, relationships, or predictive
influences (Slide 19).

\subsubsection{Central Tendency: The Three Centres (Tendance
centrale)}\label{central-tendency-the-three-centres-tendance-centrale}

\begin{enumerate}
\def\labelenumi{\arabic{enumi}.}
\tightlist
\item
  \textbf{Arithmetic Mean (\(\bar{x}\) --- Moyenne)}:
  \[\bar{x} = \frac{1}{n} \sum_{i=1}^n x_i\] The mathematical average.
  Takes every data point into account. \textbf{Weakness (Slide 21)}:
  Completely distorted and destroyed by extreme outliers and skewness.
\item
  \textbf{Median (\(\text{Mdn}\) --- Médiane)}: The exact middle value
  (50th percentile) dividing the ordered dataset into two equal halves.
  \textbf{Strength}: Robust and immune to extreme outliers.
\item
  \textbf{Mode (Mode)}: The most frequently occurring value or category
  in a distribution. \textbf{The only valid measure of central tendency
  for nominal data} (Slide 20).
\end{enumerate}

\subsubsection{Dispersion \& Variability (Mesures de
dispersion)}\label{dispersion-variability-mesures-de-dispersion}

\begin{itemize}
\tightlist
\item
  \textbf{Standard Deviation (\(s\) --- Écart-type)}:
  \[s = \sqrt{\frac{\sum (x_i - \bar{x})^2}{n - 1}}\] The average
  distance of an observation from the mean. Tells you whether the mean
  accurately represents the typical case (Slide 22).
\item
  \textbf{Variance (\(s^2\))}: The square of the standard deviation.
\item
  \textbf{Range (Étendue)}: The absolute distance between maximum and
  minimum (\(\text{Max} - \text{Min}\)).
\item
  \textbf{Interquartile Range (\(IQR\) --- Écart interquartile)}:
  \[IQR = Q_3 - Q_1\] The range covering the central 50\% of
  observations. Highly robust against outliers.
\end{itemize}

\subsubsection{Distribution Shape: Skewness \& Kurtosis (Forme de
distribution)}\label{distribution-shape-skewness-kurtosis-forme-de-distribution}

\begin{itemize}
\tightlist
\item
  \textbf{Skewness (\(g_1\) --- Asymétrie)}:
  \[g_1 = \frac{n}{(n-1)(n-2)} \sum_{i=1}^n \left(\frac{x_i - \bar{x}}{s}\right)^3\]
  Measures the lack of symmetry in a distribution.

  \begin{itemize}
  \tightlist
  \item
    \(g_1 \approx 0\): Symmetric distribution
    (\(\text{Mean} \approx \text{Median} \approx \text{Mode}\)).
  \item
    \(g_1 > 0\) (\textbf{Positive / Right Skew}): Mean \textgreater{}
    Median \textgreater{} Mode. A long tail extending to the right
    (typical of income, revenue, market cap).
  \item
    \(g_1 < 0\) (\textbf{Negative / Left Skew}): Mean \textless{} Median
    \textless{} Mode. A long tail extending to the left (typical of
    profit margins with distressed outliers).
  \end{itemize}
\item
  \textbf{Kurtosis (\(g_2\) --- Aplatissement / Acuité)}:
  \[g_2 = \frac{n(n+1)}{(n-1)(n-2)(n-3)} \sum_{i=1}^n \left(\frac{x_i - \bar{x}}{s}\right)^4 - \frac{3(n-1)^2}{(n-2)(n-3)}\]
  Measures the ``tailedness'' and peakedness of the distribution
  relative to a normal curve.

  \begin{itemize}
  \tightlist
  \item
    \(g_2 = 0\) (Mesokurtic): Standard bell-shaped normal curve.
  \item
    \(g_2 > 0\) (Leptokurtic): Sharply peaked with heavy, fat tails
    (higher probability of extreme outliers).
  \item
    \(g_2 < 0\) (Platykurtic): Flatter peak with thin tails.
  \end{itemize}
\item
  \textbf{The Normality Rule of Thumb (Slide 23)}: Both skewness and
  kurtosis falling within \textbf{\([-1.0, +1.0]\)} is consistent with
  normality.
\end{itemize}

\subsubsection{Normality Testing: The Three Converging Pieces of
Evidence (Slide
24)}\label{normality-testing-the-three-converging-pieces-of-evidence-slide-24}

Never rely on a single metric; evaluate all three simultaneously:

\begin{enumerate}
\def\labelenumi{\arabic{enumi}.}
\tightlist
\item
  \textbf{Compare the Three Centres}: If
  \(\text{Mean} \approx \text{Median} \approx \text{Mode}\), the
  distribution is roughly symmetric.
\item
  \textbf{Evaluate Skewness and Kurtosis}: Verify whether both reside
  inside the \([-1.0, +1.0]\) window.
\item
  \textbf{Visual Inspection \& Formal Shapiro-Wilk Test}:

  \begin{itemize}
  \tightlist
  \item
    \emph{Visual Inspection}: Inspect Histogram, Density curve, and Q-Q
    Plot.
  \item
    \emph{Shapiro-Wilk Test}: Formal statistical test with null
    hypothesis \(H_0\): \emph{``The data are normally distributed''}.
  \item
    \textbf{The Inverted Logic (Slide 24)}: Unlike almost every other
    test in the course, \textbf{\(p > 0.05\) means you KEEP normality}.
    If \(p < 0.05\), you \textbf{REJECT normality}.
  \end{itemize}
\end{enumerate}

\begin{center}\rule{0.5\linewidth}{0.5pt}\end{center}

\subsection{5. Session 2: Bivariate Associations, Chi-Square \&
Correlations}\label{session-2-bivariate-associations-chi-square-correlations}

\subsubsection{Bivariate Analysis (Analyse
bivariée)}\label{bivariate-analysis-analyse-bivariuxe9e}

The statistical examination of the relationship or association between
\textbf{two variables simultaneously} (Slide 17 \& Session 2).

\subsubsection{The Four-Step Protocol (Slide 9 \&
41)}\label{the-four-step-protocol-slide-9-41}

The mandatory answer template for every analytical question:

\begin{enumerate}
\def\labelenumi{\arabic{enumi}.}
\tightlist
\item
  \textbf{IDENTIFY}: Which statistical test to execute? Fully justified
  by the variable scales.
\item
  \textbf{HYPOTHESISE}: Formulate \(H_0\) and \(H_1\) in clear,
  non-technical business language.
\item
  \textbf{DECIDE}: Is \(p < 0.05\)? If yes, what is the
  magnitude/strength of the effect?
\item
  \textbf{RECOMMEND}: Actionable executive guidance answering:
  \emph{``What does the manager do on Monday morning?''}
\end{enumerate}

\subsubsection{\texorpdfstring{Hypotheses: Null (\(H_0\))
vs.~Alternative
(\(H_1\))}{Hypotheses: Null (H\_0) vs.~Alternative (H\_1)}}\label{hypotheses-null-h_0-vs.-alternative-h_1}

\begin{itemize}
\tightlist
\item
  \textbf{Null Hypothesis (\(H_0\))}: The hypothesis of ``no
  difference'', ``no relationship'', or ``status quo'' (e.g.,
  \emph{``Governance risk has no relationship with sector''}).
\item
  \textbf{Alternative Hypothesis (\(H_1\))}: The research hypothesis
  proposing an active relationship or group difference (e.g.,
  \emph{``Governance risk varies significantly across sectors''}).
\end{itemize}

\subsubsection{\texorpdfstring{Statistical Significance (\(p\)-value)
vs.~Effect Size (Slide
42)}{Statistical Significance (p-value) vs.~Effect Size (Slide 42)}}\label{statistical-significance-p-value-vs.-effect-size-slide-42}

\begin{itemize}
\tightlist
\item
  \textbf{\(p\)-value (Significance)}: The probability of observing
  results at least as extreme as the sample data assuming \(H_0\) is
  true. If \(p < 0.05\), we reject \(H_0\) and conclude that an effect
  exists.
\item
  \textbf{Reporting Standard (Slide 42)}: Never write \(p = 0.000\)!
  Software printing \texttt{.000} means the value is below three decimal
  places. Report strictly as \textbf{\(p < .001\)}.
\item
  \textbf{Effect Size (Strength)}: While \(p\)-value tells you
  \emph{whether} an effect exists, effect size tells you \emph{how
  strong} it is (whether it practically matters to a manager).
\end{itemize}

\subsubsection{\texorpdfstring{Decision Errors: Type I (\(\alpha\))
vs.~Type II (\(\beta\)) \& Statistical
Power}{Decision Errors: Type I (\textbackslash alpha) vs.~Type II (\textbackslash beta) \& Statistical Power}}\label{decision-errors-type-i-alpha-vs.-type-ii-beta-statistical-power}

In hypothesis testing, two types of erroneous conclusions are possible:

{\def\LTcaptype{none} % do not increment counter
\begin{longtable}[]{@{}
  >{\raggedright\arraybackslash}p{(\linewidth - 4\tabcolsep) * \real{0.2349}}
  >{\raggedright\arraybackslash}p{(\linewidth - 4\tabcolsep) * \real{0.3916}}
  >{\raggedright\arraybackslash}p{(\linewidth - 4\tabcolsep) * \real{0.3735}}@{}}
\toprule\noalign{}
\begin{minipage}[b]{\linewidth}\raggedright
Reality ~Decision
\end{minipage} & \begin{minipage}[b]{\linewidth}\raggedright
Fail to Reject \(H_0\) (Accept Status Quo)
\end{minipage} & \begin{minipage}[b]{\linewidth}\raggedright
Reject \(H_0\) (Conclude Effect Exists)
\end{minipage} \\
\midrule\noalign{}
\endhead
\bottomrule\noalign{}
\endlastfoot
\textbf{\(H_0\) is True} (No real effect) & \textbf{Correct Decision}
(\(1 - \alpha = 95\%\)) & \textbf{Type I Error (\(\alpha = 5\%\))}
\emph{(False Positive)} \\
\textbf{\(H_0\) is False} (Real effect exists) & \textbf{Type II Error
(\(\beta\))} \emph{(False Negative / Missed Discovery)} &
\textbf{Correct Decision: Statistical Power (\(1 - \beta \ge 80\%\))} \\
\end{longtable}
}

\begin{itemize}
\tightlist
\item
  \textbf{Type I Error (\(\alpha = 0.05\) --- Faux positif)}: Claiming a
  governance relationship exists when it is purely due to random chance.
  The significance threshold \(\alpha\) fixes this maximum acceptable
  risk at 5\%.
\item
  \textbf{Type II Error (\(\beta\) --- Faux négatif)}: Missing a genuine
  governance impact because the sample is too small or noisy.
\item
  \textbf{Statistical Power (\(1 - \beta\) --- Puissance statistique)}:
  The probability of successfully detecting a genuine effect. By having
  \(N = 503\) (well above the course minimum of 100), our study achieves
  high statistical power (\(> 95\%\)), minimizing Type II errors.
\end{itemize}

\subsubsection{95\% Confidence Interval (Intervalle de confiance à 95\%
--- CI)}\label{confidence-interval-intervalle-de-confiance-uxe0-95-ci}

\[\text{CI}_{95\%} = \text{Point Estimate} \pm t_{\text{crit}} \times \text{Standard Error}\]
A range of values calculated from the sample data that has a 95\%
probability of containing the true population parameter:

\begin{itemize}
\tightlist
\item
  \textbf{Interpretation Rule}: If the 95\% Confidence Interval for a
  difference or correlation contains zero (e.g., \([-0.05, +0.12]\)),
  the effect is \textbf{not statistically significant} at
  \(\alpha = 0.05\). If zero is excluded (e.g., Pearson \(r\) for
  Governance Risk and Beta: \([+0.18, +0.35]\)), the relationship is
  statistically significant.
\item
  \textbf{Margin of Error (Marge d'erreur)}: Half the width of the
  confidence interval (\(t_{\text{crit}} \times SE\)), quantifying
  sampling uncertainty.
\end{itemize}

\subsubsection{\texorpdfstring{Chi-Square Test of Independence
(\(\chi^2\) --- Nominal \(\times\)
Ordinal)}{Chi-Square Test of Independence (\textbackslash chi\^{}2 --- Nominal \textbackslash times Ordinal)}}\label{chi-square-test-of-independence-chi2-nominal-times-ordinal}

Tests whether two qualitative categorical variables are independent or
statistically associated.
\[\chi^2 = \sum_{i=1}^r \sum_{j=1}^c \frac{(O_{ij} - E_{ij})^2}{E_{ij}}\]
Where \(O_{ij}\) is the observed cell count, and
\(E_{ij} = \frac{\text{Row Total} \times \text{Column Total}}{N}\) is
the expected count under independence. Degrees of freedom:
\(df = (r - 1)(c - 1)\).

\begin{itemize}
\tightlist
\item
  \textbf{Cramér's \(V\) (Effect Size for Chi-Square)}:
  \[V = \sqrt{\frac{\chi^2}{N \min(r - 1, c - 1)}}\] Measures
  association strength from \(0\) (no relationship) to \(1\) (perfect
  association):

  \begin{itemize}
  \tightlist
  \item
    \(V < 0.10\): Negligible
  \item
    \(0.10 \le V < 0.30\): Weak to moderate association (Our result:
    \(V = 0.122\))
  \item
    \(V \ge 0.30\): Strong association
  \end{itemize}
\item
  \textbf{Standardized Residuals (Résidus standardisés)}:
  \[\text{Residual}_{ij} = \frac{O_{ij} - E_{ij}}{\sqrt{E_{ij}}}\]
  Identifies which specific cells drive the overall Chi-Square result. A
  residual \(> +2.0\) indicates significant over-representation;
  \(< -2.0\) indicates significant under-representation. (In our data,
  Tech \& Comms had \(+1.75\) in High Risk; Industrials had \(-1.97\)).
\end{itemize}

\subsubsection{\texorpdfstring{Pearson Correlation Coefficient (\(r\)
--- Continuous \(\times\)
Continuous)}{Pearson Correlation Coefficient (r --- Continuous \textbackslash times Continuous)}}\label{pearson-correlation-coefficient-r-continuous-times-continuous}

\[r = \frac{\sum (x_i - \bar{x})(y_i - \bar{y})}{\sqrt{\sum (x_i - \bar{x})^2 \sum (y_i - \bar{y})^2}}\]
Measures the linear strength and direction of the association between
two continuous variables:

\begin{itemize}
\tightlist
\item
  Ranges from \(-1.0\) (perfect inverse link) to \(+1.0\) (perfect
  direct link).
\item
  \textbf{Fisher \(z\)-Transformation \& 95\% Confidence Interval}:
  Constructs exact lower and upper confidence bounds:
  \[z = \frac{1}{2} \ln \left(\frac{1+r}{1-r}\right), \quad SE_z = \frac{1}{\sqrt{n - 3}}\]
\end{itemize}

\subsubsection{\texorpdfstring{Spearman's Rank Correlation (\(\rho\) ---
Non-Paramétrique)}{Spearman's Rank Correlation (\textbackslash rho --- Non-Paramétrique)}}\label{spearmans-rank-correlation-rho-non-paramuxe9trique}

Calculates Pearson's correlation on the \textbf{ranks} of data rather
than raw values. Robust to non-normality and monotonic non-linear
patterns.

\begin{center}\rule{0.5\linewidth}{0.5pt}\end{center}

\subsection{6. Session 3: Group Comparisons, Welch's ANOVA \&
Games-Howell}\label{session-3-group-comparisons-welchs-anova-games-howell}

\subsubsection{\texorpdfstring{One-Way Analysis of Variance (ANOVA ---
Qualitative 3+ groups \(\times\)
Quantitative)}{One-Way Analysis of Variance (ANOVA --- Qualitative 3+ groups \textbackslash times Quantitative)}}\label{one-way-analysis-of-variance-anova-qualitative-3-groups-times-quantitative}

Tests whether the population means of a continuous outcome variable
differ significantly across 3 or more categorical groups.
\[F = \frac{MS_{\text{between}}}{MS_{\text{within}}} = \frac{SS_{\text{between}} / (k - 1)}{SS_{\text{within}} / (N - k)}\]

\subsubsection{Levene's Test of Homogeneity of Variances (Test de
Levene)}\label{levenes-test-of-homogeneity-of-variances-test-de-levene}

Tests whether variances are equal across groups
(\(H_0: \sigma_1^2 = \sigma_2^2 = \dots = \sigma_k^2\)).

\begin{itemize}
\tightlist
\item
  \emph{If \(p > 0.05\)}: Variances are equal (homoscedasticity
  confirmed; standard Fisher ANOVA is valid).
\item
  \emph{If \(p < 0.05\)}: Variances are significantly unequal
  (\textbf{heteroscedasticity}; standard Fisher ANOVA is invalid and
  Welch's ANOVA must be used!).
\item
  \emph{In our project}: Levene's test yielded \(F = 5.39, p < .001\)
  for Profit Margin across sectors.
\end{itemize}

\subsubsection{\texorpdfstring{Welch's Robust ANOVA
(\(F_{\text{Welch}}\))}{Welch's Robust ANOVA (F\_\{\textbackslash text\{Welch\}\})}}\label{welchs-robust-anova-f_textwelch}

An adjusted ANOVA formulation that weights each group's variance by its
sample size (\(w_i = n_i / s_i^2\)). Valid even when group variances and
sample sizes are severely unequal.

\subsubsection{\texorpdfstring{Effect Size: Eta-Squared (\(\eta^2\)) \&
Omega-Squared
(\(\omega^2\))}{Effect Size: Eta-Squared (\textbackslash eta\^{}2) \& Omega-Squared (\textbackslash omega\^{}2)}}\label{effect-size-eta-squared-eta2-omega-squared-omega2}

\[\eta^2 = \frac{SS_{\text{between}}}{SS_{\text{total}}}\] The
percentage of variance in the dependent variable accounted for by group
membership:

\begin{itemize}
\tightlist
\item
  \(\eta^2 \approx 0.01\): Small effect
\item
  \(\eta^2 \approx 0.06\): Medium effect
\item
  \(\eta^2 \ge 0.14\): Large effect
\item
  \emph{In our project}: \(\eta^2 = 0.0405\) (\(4.05\%\) of profit
  margin variance is explained by sector).
\end{itemize}

\subsubsection{\texorpdfstring{Kruskal-Wallis Non-Parametric Test
(\(H\))}{Kruskal-Wallis Non-Parametric Test (H)}}\label{kruskal-wallis-non-parametric-test-h}

The non-parametric rank-based equivalent of One-Way ANOVA. Used when
both normality and variance homogeneity are severely violated.

\subsubsection{Post-Hoc Pairwise Comparisons: Tukey HSD
vs.~Games-Howell}\label{post-hoc-pairwise-comparisons-tukey-hsd-vs.-games-howell}

When ANOVA indicates that at least one group differs (\(p < 0.05\)),
post-hoc tests identify \emph{which specific pairs} differ:

\begin{itemize}
\tightlist
\item
  \textbf{Tukey's HSD (Honestly Significant Difference)}: Assumes equal
  variances across all groups.
\item
  \textbf{Games-Howell Post-Hoc Test}:
  \[t = \frac{\bar{x}_i - \bar{x}_j}{\sqrt{\frac{s_i^2}{n_i} + \frac{s_j^2}{n_j}}}\]
  Adjusts degrees of freedom via the Welch-Satterthwaite equation and
  applies the Studentized Range distribution (\(q\)). \textbf{Required
  when Levene's test fails!}
\end{itemize}

\begin{center}\rule{0.5\linewidth}{0.5pt}\end{center}

\subsection{7. Session 4: Multiple Linear Regression \& Econometric
Diagnostics}\label{session-4-multiple-linear-regression-econometric-diagnostics}

\subsubsection{Multiple Linear Regression (Régression linéaire multiple
MCO /
OLS)}\label{multiple-linear-regression-ruxe9gression-linuxe9aire-multiple-mco-ols}

Models the expected value of a continuous dependent variable (\(Y\)) as
a linear combination of explanatory variables
(\(X_1, X_2, \dots, X_k\)):
\[Y_i = \beta_0 + \beta_1 X_{1i} + \beta_2 X_{2i} + \dots + \beta_k X_{ki} + \epsilon_i\]

\begin{itemize}
\tightlist
\item
  \textbf{\(\beta_0\) (Intercept / Constante)}: The predicted baseline
  value of \(Y\) when all predictors are equal to zero.
\item
  \textbf{\(\beta_k\) (Unstandardized Regression Coefficient)}: The
  expected change in \(Y\) for a 1-unit increase in \(X_k\), holding all
  other variables constant (\emph{ceteris paribus}).
\item
  \textbf{Standardized Beta (\(\beta^*\))}:
  \[\beta_k^* = \beta_k \times \frac{s_{X_k}}{s_Y}\] Allows direct
  comparison of predictor importance on a common standard-deviation
  scale.
\end{itemize}

\subsubsection{\texorpdfstring{Coefficient of Determination (\(R^2\) and
Adjusted
\(R^2\))}{Coefficient of Determination (R\^{}2 and Adjusted R\^{}2)}}\label{coefficient-of-determination-r2-and-adjusted-r2}

\begin{itemize}
\tightlist
\item
  \textbf{\(R^2\)}: The proportion of variance in \(Y\) explained by the
  regression model:
  \[R^2 = 1 - \frac{SS_{\text{residual}}}{SS_{\text{total}}}\]
\item
  \textbf{Adjusted \(R^2\)}: Corrects \(R^2\) for the number of
  predictors (\(k\)) and sample size (\(N\)):
  \[R_{\text{adj}}^2 = 1 - \left[\frac{(1 - R^2)(N - 1)}{N - k - 1}\right]\]
\end{itemize}

\subsubsection{\texorpdfstring{Logarithmic Transformations
(\(\ln(X)\))}{Logarithmic Transformations (\textbackslash ln(X))}}\label{logarithmic-transformations-lnx}

Applying natural logarithms to positively skewed financial metrics
(Market Capitalization and Revenue).

\begin{itemize}
\tightlist
\item
  \emph{Econometric Rationale}: Compresses extreme mega-cap outliers
  (NVIDIA, Apple) and transforms multiplicative scale dynamics into
  linear additive relationships.
\end{itemize}

\subsubsection{Econometric Assumptions \& Diagnostic Suite (Slide
133--137)}\label{econometric-assumptions-diagnostic-suite-slide-133137}

\begin{enumerate}
\def\labelenumi{\arabic{enumi}.}
\tightlist
\item
  \textbf{Linearity}: Evaluated via the Residuals vs.~Fitted plot
  (should show a flat horizontal band with no curvature).
\item
  \textbf{Normality of Residuals}: Model errors (\(\epsilon\)) must be
  normally distributed (Q-Q plot of residuals).
\item
  \textbf{Homoscedasticity (Equal Error Variance)}:

  \begin{itemize}
  \tightlist
  \item
    \textbf{Breusch-Pagan Test}: Regresses squared residuals on
    explanatory variables.
  \item
    \emph{If \(p < 0.05\)}: Heteroscedasticity is present (White's
    robust standard errors HC3 are required).
  \end{itemize}
\item
  \textbf{Absence of Multicollinearity (VIF \& Tolerance)}:

  \begin{itemize}
  \tightlist
  \item
    \textbf{Variance Inflation Factor (VIF)}:
    \[\text{VIF}_j = \frac{1}{1 - R_j^2}\] Measures how much the
    variance of an estimated regression coefficient increases due to
    collinearity with other predictors:

    \begin{itemize}
    \tightlist
    \item
      \(\text{VIF} < 5.0\): Safe / Low collinearity (Our maximum VIF was
      \(1.89\)).
    \item
      \(\text{VIF} > 10.0\): Severe multicollinearity (distorts
      \(t\)-statistics and inflates standard errors).
    \end{itemize}
  \item
    \textbf{Tolerance}: The inverse of VIF
    (\(\text{Tol} = 1 / \text{VIF}\)). Tolerance \(< 0.20\) signals high
    collinearity risk.
  \end{itemize}
\end{enumerate}

\subsubsection{Dummy Variables \& Reference Category (Variables
indicatrices /
muettes)}\label{dummy-variables-reference-category-variables-indicatrices-muettes}

To include a qualitative nominal variable (like \texttt{Sector} with 5
groups) in an OLS regression, it must be converted into
\textbf{\(k - 1 = 4\) binary dummy variables} taking values \(0\) or
\(1\):

\begin{itemize}
\tightlist
\item
  \textbf{Reference Group (Catégorie de référence)}: \emph{Tech \&
  Comms} is omitted from the equation to serve as the baseline
  (\(X_{\text{all\_dummies}} = 0\)).
\item
  \textbf{Interpretation of Dummy Coefficients
  (\(\beta_{\text{Sector}}\))}: The coefficient represents the expected
  difference in profit margin between that sector and the reference
  sector (Tech \& Comms), holding all financial and governance variables
  constant.
\item
  \textbf{Dummy Variable Trap (Piège de la colinéarité parfaite)}: If
  all 5 dummy variables were included alongside the constant
  (\(\beta_0\)), the sum of dummies would equal \(1\), causing perfect
  multicollinearity (\(\text{VIF} = \infty\)) and rendering matrix
  inversion mathematically impossible.
\end{itemize}

\subsubsection{\texorpdfstring{Overall Model Utility: The Regression
\(F\)-Test}{Overall Model Utility: The Regression F-Test}}\label{overall-model-utility-the-regression-f-test}

Tests the global null hypothesis that \emph{none} of the explanatory
variables predict \(Y\)
(\(H_0: \beta_1 = \beta_2 = \dots = \beta_k = 0\)):
\[F = \frac{MS_{\text{model}}}{MS_{\text{residual}}} = \frac{R^2 / k}{(1 - R^2) / (N - k - 1)}\]
If \(p_F < 0.05\), we reject \(H_0\) and conclude that the model as a
whole has genuine explanatory power beyond pure chance. (In our Model 3:
\(F(8, 485) = 11.23, p < .001^{***}\)).

\subsubsection{\texorpdfstring{Influential Observations: Leverage
(\(h_{ii}\)) \& Cook's Distance
(\(D_i\))}{Influential Observations: Leverage (h\_\{ii\}) \& Cook's Distance (D\_i)}}\label{influential-observations-leverage-h_ii-cooks-distance-d_i}

Not all outliers exert the same distortive pull on regression slopes:

\begin{itemize}
\tightlist
\item
  \textbf{Leverage (\(h_{ii}\) --- Valeur levier)}: Measures how far an
  observation's predictor values (\(X\)) are from the center of the
  predictor space. High leverage points (e.g., massive revenue
  conglomerates) have high potential to swing the regression line.
  Warning threshold: \(h_{ii} > 2(k + 1) / N\).
\item
  \textbf{Studentized Residual (\(r_i\))}: The residual divided by its
  estimated standard deviation, flagging outliers in the \(Y\)-dimension
  (\(|r_i| > 3.0\)).
\item
  \textbf{Cook's Distance (\(D_i\) --- Distance de Cook)}:
  \[D_i = \frac{\sum (\hat{y}_j - \hat{y}_{j(i)})^2}{(k + 1) s^2}\]
  Measures the aggregate shift in all model predictions when observation
  \(i\) is excluded. Any observation with \(D_i > 1.0\) (or
  \(D_i > 4/N\)) is an \textbf{influential observation} requiring close
  inspection. (Panel D of Figure 9 visualizes studentized residuals
  against leverage).
\end{itemize}

\subsubsection{\texorpdfstring{Autocorrelation of Residuals \&
Durbin-Watson Test
(\(d\))}{Autocorrelation of Residuals \& Durbin-Watson Test (d)}}\label{autocorrelation-of-residuals-durbin-watson-test-d}

Assumption that error terms are independent
(\(\text{Cov}(\epsilon_i, \epsilon_j) = 0\)):

\begin{itemize}
\tightlist
\item
  \textbf{Durbin-Watson Statistic (\(d\))}:
  \[d = \frac{\sum_{i=2}^N (e_i - e_{i-1})^2}{\sum_{i=1}^N e_i^2}\]
  Ranges from \(0\) to \(4\):

  \begin{itemize}
  \tightlist
  \item
    \(d \approx 2.0\): No autocorrelation (residuals are completely
    independent).
  \item
    \(d < 1.5\): Positive autocorrelation (common in time-series data).
  \item
    \(d > 2.5\): Negative autocorrelation.
  \end{itemize}
\item
  \textbf{In our Cross-Sectional Data}: Because 1 row = 1 independent
  firm, autocorrelation is naturally absent (\(d \approx 1.98\)).
\end{itemize}

\begin{center}\rule{0.5\linewidth}{0.5pt}\end{center}

\subsection{8. Jamovi Software Operations \& Costly Student
Traps}\label{jamovi-software-operations-costly-student-traps}

\subsubsection{What is Jamovi?}\label{what-is-jamovi}

A free, open-source, point-and-click statistical software package built
on top of the R statistical programming language, providing SPSS-style
output tables and real-time reactive calculations (Slide 37).

\subsubsection{Jamovi Measure Types (Slide
38)}\label{jamovi-measure-types-slide-38}

\begin{itemize}
\tightlist
\item
  \textbf{ID}: Column used solely to identify entities (e.g.,
  \texttt{Ticker}). Jamovi never computes statistics on ID columns.
\item
  \textbf{Nominal}: Unordered qualitative categories (icon: three
  colored circles).
\item
  \textbf{Ordinal}: Ranked qualitative categories (icon: small step
  chart).
\item
  \textbf{Continuous}: Numeric measurements with decimals (icon: small
  ruler).
\end{itemize}

\subsubsection{\texorpdfstring{The \texttt{.omv} File
Format}{The .omv File Format}}\label{the-.omv-file-format}

The proprietary Jamovi file format that bundles the spreadsheet dataset,
variable configurations, and all statistical output tables and graphs
together into a single file. \textbf{This is the primary file submitted
for lab coursework} (Slide 37).

\subsubsection{Top 6 Errors That Cost Marks Every Year (Slide
42)}\label{top-6-errors-that-cost-marks-every-year-slide-42}

\begin{enumerate}
\def\labelenumi{\arabic{enumi}.}
\tightlist
\item
  \textbf{Reporting a mean for a nominal variable}: e.g., writing
  \emph{``average sector = 2.4''}.
\item
  \textbf{Setting the wrong measure type in Jamovi}: e.g., setting
  \texttt{Sector} as Continuous or \texttt{Profit\_Margin} as Nominal,
  causing Jamovi to hide the required statistical tests.
\item
  \textbf{Reading ``Sig. = 0.000'' as \(p = 0\)}: Always format as
  \textbf{\(p < .001\)}.
\item
  \textbf{Confusing significance with strength}: A test can be
  statistically significant (\(p < 0.05\)) because sample size \(N\) is
  large, but have a tiny, negligible effect size (\(r = 0.08\)).
\item
  \textbf{Pasting output tables without interpretation}: Half the marks
  are awarded for explaining what the numbers mean in plain English.
\item
  \textbf{Stopping at the statistic without recommending action}: Always
  conclude with what the executive or portfolio manager does on Monday
  morning.
\end{enumerate}

\end{document}
